\documentclass[11pt]{article}
\usepackage[margin=1in]{geometry}
\usepackage{amsmath,amssymb}
\usepackage{graphicx}
\usepackage{hyperref}
\usepackage{enumitem}
\usepackage{booktabs}
\usepackage{float}
\usepackage[font=small,labelfont=bf]{caption}

\title{ChemCS: Integrating Constraint Satisfaction and Diffusion\\Models for Chemically Valid Molecular Generation}
\author{Sreepranav Chetlapalli}
\date{}

\begin{document}
\maketitle

% ===================================================================
% ABSTRACT
% ===================================================================
\begin{abstract}
Molecular generation with diffusion models has shown promise for drug
discovery and materials design, yet these models frequently produce
chemically invalid structures such as atoms with too many bonds,
reactions that violate mass conservation, or mismatched formal charges.
Most existing approaches address this by filtering invalid outputs
after generation is complete, which wastes computation and provides no
mechanism for mid-trajectory correction.

This paper presents ChemCS, a system that integrates a
constraint-satisfaction supervisor directly into the reverse-diffusion
loop of a graph neural network (GNN) molecular generator. At every
denoising step the supervisor decodes the current molecular state,
verifies chemical axioms (valency limits, mass conservation, charge
balance, and atom counts), and either commits the step, applies targeted
repairs, or backtracks to a previously valid state. The constraint
checker uses the Z3 SMT solver from Microsoft Research to formulate and
verify chemical rules as satisfiability queries.

The generator itself is a lightweight two-layer GNN that operates on
atom-level feature vectors and an adjacency matrix of bond orders. A
companion PyTorch training pipeline adds differentiable constraint
penalties for valency and element conservation to the standard denoising
loss, with a curriculum schedule that ramps these penalties gradually to
avoid destabilizing early training.

On a benchmark of three representative reactions (SN2, combustion,
acid-base) with 50 random seeds each, supervised generation achieves
100\% valency validity compared with roughly 15\% for unsupervised
generation with the same untrained network.
\end{abstract}

% ===================================================================
% 1  INTRODUCTION
% ===================================================================
\section{Introduction}

Generating novel molecular structures is a central challenge in
computational chemistry. Applications range from \emph{de novo} drug
design, where candidate molecules must satisfy pharmacological
constraints, to materials science, where target properties such as band
gap or thermal stability guide the search. In recent years, diffusion
models, originally developed for image generation, have been adapted to
molecular graphs with encouraging
results~\cite{hoogeboom2022equivariant, vignac2023digress}.

A diffusion model generates a sample by starting from pure noise and
iteratively refining it through a sequence of learned denoising steps.
When applied to molecules the noise is added to atom features and bond
orders, and the reverse process gradually recovers a structured molecular
graph. The difficulty is that the neural network performing these updates
has no built-in knowledge of chemistry: it can propose carbon atoms with
five bonds, reactions where mass appears or disappears, or products
whose total charge does not match the reactants.

The standard mitigation is post-hoc filtering: generate many candidate
molecules, discard the invalid ones, and keep what remains. This
strategy has two significant drawbacks. First, it is wasteful. If the
model has a low validity rate, the vast majority of computation is
thrown away. Second, and more importantly, it misses the opportunity to
guide the generative process. If an error is detected at step 40 of a
50-step trajectory, there is no reason to wait until step 50 to discover
and discard it.

ChemCS addresses this gap by wrapping the diffusion model in a
supervisor loop that enforces chemical axioms at every reverse step. The
key contributions of this work are:

\begin{enumerate}[nosep]
  \item A constraint-satisfaction module that encodes four fundamental
        chemical rules (valency limits, mass conservation, charge
        balance, and atom-count conservation) as queries to the Z3
        satisfiability modulo theories (SMT) solver.
  \item A supervisor loop that intercepts each denoising step, runs
        constraint checks, and responds with one of three actions:
        \emph{commit} (accept the step), \emph{correct} (apply targeted
        repairs), or \emph{backtrack} (revert to the last valid state).
  \item A constraint-aware training objective that supplements the
        standard denoising loss with differentiable penalties for valency
        violations and element-count mismatches, introduced through a
        curriculum schedule.
  \item A comprehensive test suite and benchmark framework for
        reproducible evaluation.
\end{enumerate}

% ===================================================================
% 2  BACKGROUND
% ===================================================================
\section{Background}

This section provides a brief overview of the two main technical
foundations that ChemCS builds upon: denoising diffusion models for
graph-structured data and satisfiability modulo theories for constraint
checking.

\subsection{Denoising Diffusion Models on Graphs}

A denoising diffusion probabilistic model (DDPM) defines a forward
process that progressively corrupts a data sample $x_0$ by adding
Gaussian noise over $T$ timesteps, producing a sequence
$x_0, x_1, \ldots, x_T$ where $x_T$ is approximately standard Gaussian
noise. Generation works by learning the reverse process: a neural
network $\varepsilon_\theta$ is trained to predict the noise added at
each step, so that starting from $x_T \sim \mathcal{N}(0, I)$ the model
can iteratively recover a clean sample $x_0$.

The amount of noise at each timestep is controlled by a variance
schedule $\{\beta_t\}_{t=1}^T$. A key quantity is the cumulative
signal-retention factor
$\bar\alpha_t = \prod_{s=1}^t (1 - \beta_s)$, which determines how much
of the original signal remains at step $t$. When $\bar\alpha_T \approx 0$
the data is fully destroyed; when $\bar\alpha_1 \approx 1$ almost no
noise has been added.

For molecular graphs, the data $x_0$ consists of two components: a
matrix of atom features (element type, bond count, charge) and an
adjacency matrix of bond orders. The forward process adds continuous
Gaussian noise to atom features and randomly corrupts bond orders. The
reverse process uses a graph neural network to predict the clean state
from the noisy one. ChemCS supports both linear and
cosine~\cite{nichol2021improved} noise schedules.

\subsection{Satisfiability Modulo Theories (SMT) and Z3}

Satisfiability modulo theories (SMT) extends Boolean satisfiability
(SAT) with reasoning over richer domains such as integers, real
numbers, and arrays. An SMT solver takes a set of logical constraints
and determines whether any variable assignment satisfies all of them
simultaneously.

Z3~\cite{demoura2008z3} is a high-performance SMT solver developed by
Microsoft Research. In ChemCS, Z3 is used to encode chemical rules as
constraints: for example, the statement ``for every atom $i$, the total
number of bonds must not exceed the allowed valency for that element and
charge'' becomes an integer-arithmetic constraint that Z3 can check or
refute in microseconds.

The advantage of using an SMT solver rather than simple procedural
checks is extensibility: adding a new chemical rule amounts to adding a
new constraint to the solver, without modifying control flow. Although
the current set of rules is small enough that an SMT solver is not
strictly necessary, it establishes a framework that can accommodate more
complex constraints (such as ring-size limits or aromaticity rules) in
the future.

% ===================================================================
% 3  PROBLEM FORMULATION
% ===================================================================
\section{Problem Formulation}

\subsection{Molecular Representation}

A molecule is represented as a set of atoms
$\{a_1, a_2, \ldots, a_n\}$, where each atom $a_i$ carries four
attributes:

\begin{itemize}[nosep]
  \item Element $e_i$: the chemical element (e.g., C, N, O).
  \item Bond count $b_i$: the sum of bond orders to neighbouring atoms.
  \item Formal charge $q_i$: the integer formal charge
        (e.g., $-1$ for bromide).
  \item Implicit hydrogen count $h_i$: hydrogens not explicitly
        represented in the graph.
\end{itemize}

Bonds between atoms are stored in a symmetric adjacency matrix
$A \in \{0,1,2,3\}^{n \times n}$, where $A_{ij}$ is the bond order
between atoms $i$ and $j$ (0 = no bond, 1 = single, 2 = double,
3 = triple).

\subsection{Chemical Axioms}

The supervisor enforces four axioms. These are fundamental conservation
laws and bonding rules from general chemistry:

\begin{table}[H]
\centering
\caption{Chemical axioms enforced by the constraint checker.}
\label{tab:axioms}
\begin{tabular}{@{}llp{6.5cm}@{}}
\toprule
Axiom & Formal rule & Meaning \\
\midrule
Valency &
  $b_i + h_i \leq V(e_i, q_i)$ &
  Every atom's total bond count (explicit bonds plus implicit hydrogens)
  must not exceed the maximum valency allowed for its element and charge.
  For example, neutral carbon allows 4; nitrogen with charge $+1$ allows 4. \\[4pt]
Mass conservation &
  $\bigl|\sum_{\text{react}} m - \sum_{\text{prod}} m\bigr| \leq \epsilon$ &
  The total atomic mass of reactants must equal the total atomic mass
  of products, within a small tolerance $\epsilon$ (default 0.02\,u). \\[4pt]
Charge conservation &
  $\sum_{\text{react}} q = \sum_{\text{prod}} q$ &
  The net formal charge must be identical on both sides of a reaction. \\[4pt]
Atom conservation &
  $\forall\, e:\; n_e^{\text{react}} = n_e^{\text{prod}}$ &
  Every element must appear in the same count on both sides. \\
\bottomrule
\end{tabular}
\end{table}

The effective valency function $V(e_i, q_i)$ accounts for
charge-adjusted bonding capacity. Table~\ref{tab:elements} lists the
ten elements currently supported along with their standard maximum
valencies and how formal charge can alter them.

\begin{table}[H]
\centering
\caption{Supported elements with atomic mass, maximum valency, and
         charge-adjusted valency. The ``Charge adjustment'' column shows
         how a non-zero formal charge on the atom changes the number of
         bonds it is allowed to form. For example, nitrogen normally
         allows at most 3 bonds, but when it carries a $+1$ formal
         charge (as in the ammonium ion NH$_4^+$) the limit increases
         to~4. A dash indicates that no charge adjustment is defined.}
\label{tab:elements}
\begin{tabular}{@{}lrrp{6.8cm}@{}}
\toprule
Element & Mass (u) & Max valency & Charge adjustment \\
\midrule
H   &   1.008 & 1 & None \\
C   &  12.011 & 4 & None \\
N   &  14.007 & 3 & With charge $+1$, max valency becomes 4
                     (e.g.\ NH$_4^+$). With charge $-1$, max valency
                     becomes 2. \\
O   &  15.999 & 2 & With charge $-1$, max valency becomes 1
                     (e.g.\ hydroxide OH$^-$). \\
F   &  18.998 & 1 & None \\
P   &  30.974 & 5 & With charge $+1$, max valency becomes 6. \\
S   &  32.060 & 6 & With charge $+1$, max valency becomes 7.
                     With charge $+2$, max valency becomes 8. \\
Cl  &  35.450 & 1 & None \\
Br  &  79.904 & 1 & None \\
I   & 126.904 & 1 & None \\
\bottomrule
\end{tabular}
\end{table}

\subsection{The Generation Task}

Given a set of reactant molecules, the task is to generate a set of
product molecules such that all four axioms in Table~\ref{tab:axioms}
are satisfied. The diffusion model proposes candidates through iterative
denoising, and the supervisor ensures the result is chemically
consistent.

% ===================================================================
% 4  METHODOLOGY
% ===================================================================
\section{Methodology}

\subsection{Atom Feature Encoding}

Each atom in a molecule is converted to a fixed-length feature vector
of dimension 12. The first ten entries form a one-hot encoding of the
element type; the remaining two entries encode the bond count and formal
charge, each normalized to a standard range.

\begin{table}[H]
\centering
\caption{Layout of the 12-dimensional atom feature vector.}
\label{tab:features}
\begin{tabular}{@{}clll@{}}
\toprule
Index & Feature & Encoding & Example \\
\midrule
0--9   & Element type & One-hot over 10 elements & Carbon $\to [0,1,0,\ldots,0]$ \\
10     & Bond count   & $b_i / 4$                & 4 bonds $\to 1.0$ \\
11     & Formal charge & $q_i / 2$               & $-1 \to -0.5$ \\
\bottomrule
\end{tabular}
\end{table}

The adjacency matrix $A$ is constructed by greedily assigning bond
orders between atom pairs based on each atom's remaining bonding
capacity, producing a symmetric $n \times n$ matrix. The combination of
$X \in \mathbb{R}^{n \times 12}$ (atom features) and
$A \in \mathbb{R}^{n \times n}$ (bond orders) constitutes the complete
molecular encoding.

\subsection{GNN Denoiser Architecture}

The denoiser is a two-layer graph neural network implemented in NumPy
for lightweight CPU inference (a PyTorch version with identical
architecture is used for training). Each graph convolutional layer
updates every atom's hidden representation by combining two
transformations:

\begin{equation}
h_i^{(\ell+1)} = \text{ReLU}\!\left(
  W_{\text{self}}^{(\ell)}\, h_i^{(\ell)}
  + W_{\text{neigh}}^{(\ell)}\, \frac{1}{|\mathcal{N}(i)|}
    \sum_{j \in \mathcal{N}(i)} h_j^{(\ell)}
\right)
\label{eq:gnn}
\end{equation}

\noindent
where $h_i^{(\ell)}$ is the hidden vector of atom $i$ at layer $\ell$,
$\mathcal{N}(i)$ is the set of neighbours of atom $i$ (determined by the
adjacency matrix), and $W_{\text{self}}^{(\ell)}$,
$W_{\text{neigh}}^{(\ell)}$ are learnable weight matrices. The
neighbour features are mean-aggregated (divided by the degree), which
normalizes the contribution regardless of how many bonds an atom has.

After two graph convolutional layers, the network produces outputs
through two heads:

\begin{itemize}[nosep]
  \item Atom head: A linear projection from the hidden dimension back to
        the 12-dimensional feature space, with softmax applied over the
        element logits (indices 0--9) to produce a probability
        distribution over element types.
  \item Bond head: For each pair of atoms $(i, j)$, the hidden vectors
        $h_i$ and $h_j$ are concatenated into a single vector of size
        $2 \times \text{hidden\_dim}$, then passed through a linear layer
        with 4 outputs corresponding to bond orders 0--3. The predicted
        bond order is the argmax of these logits.
\end{itemize}

The default hidden dimension is 64, giving the model approximately
20,000 trainable parameters, which is small enough to run on CPU in
under a second for molecules with up to roughly 15 atoms.

\subsection{Noise Schedule}

The forward diffusion process adds noise according to a variance
schedule. ChemCS supports two schedules:

The linear schedule sets the noise level $\beta_t$ to increase linearly
from $\beta_{\min} = 10^{-4}$ to $\beta_{\max} = 0.1$:
\begin{equation}
\beta_t = \beta_{\min} + \frac{t}{T}\,(\beta_{\max} - \beta_{\min})
\end{equation}

The cosine schedule~\cite{nichol2021improved} sets the cumulative
signal-retention factor to follow a cosine curve:
\begin{equation}
\bar\alpha_t = \frac{f(t)}{f(0)}, \quad
f(u) = \cos^2\!\left(\frac{u/T + s}{1 + s} \cdot \frac{\pi}{2}\right)
\end{equation}
with offset $s = 0.008$. The cosine schedule provides a gentler noise
ramp in the early and late timesteps, which has been shown to improve
sample quality.

For atom features, noise is added via
$x_t = \sqrt{\bar\alpha_t}\, x_0 + \sqrt{1 - \bar\alpha_t}\, \epsilon$
where $\epsilon \sim \mathcal{N}(0, I)$. For bond orders, each entry in
the adjacency matrix is independently re-sampled from $\{0,1,2,3\}$
with probability $(1 - \bar\alpha_t)$, and symmetry is enforced by
mirroring the upper triangle.

\subsection{Constraint Checking with Z3}

The constraint checker encodes each of the four chemical axioms
(Table~\ref{tab:axioms}) as a Z3 satisfiability query. The checker
works as follows:

\begin{enumerate}[nosep]
  \item Mass conservation: The mass difference between reactants and
        products is computed using Z3 real-valued arithmetic. The solver
        is asked whether the difference exceeds the
        tolerance $\epsilon$. If the formula is satisfiable, mass is
        \emph{not} conserved, and a violation is recorded.
  \item Charge conservation: The total formal charges on both sides are
        compared as integers. A mismatch is recorded directly.
  \item Atom conservation: For each element, the counts on the reactant
        and product sides are encoded as Z3 integer constants. The solver
        checks whether they differ; if satisfiable, a violation is
        recorded for that element.
  \item Valency: For each atom in the products, the solver checks
        whether the total bond count (explicit bonds plus implicit
        hydrogens) exceeds the effective valency for that atom's element
        and formal charge. Per-atom violations are reported individually.
\end{enumerate}

The function \texttt{check\_intermediate} performs only the valency
check (no reaction context needed) and is called at every denoising
step. The full four-axiom check via \texttt{check\_reaction} is invoked
at the final step ($t = 1$) and after the trajectory completes.

Both functions return a \texttt{ConstraintResult} object containing a
Boolean \texttt{sat} field and a list of human-readable violation
strings, making the results easy to log and debug.

\subsection{Supervisor Loop}

The supervisor loop is the central contribution of ChemCS. It wraps the
diffusion model's reverse process and intercepts every denoising step to
enforce chemical validity. The procedure works as follows:

\begin{enumerate}[nosep]
  \item All reactant atoms are concatenated into a single initial state.
        This state is encoded into atom features $X$ and adjacency
        matrix $A$, then noise is added up to timestep $T$.
  \item From $t = T$ down to $t = 1$, the model performs one denoising
        step to produce a candidate state. The candidate is decoded back
        into a molecular representation and checked against the relevant
        constraints (valency at every step; full mass, charge, and atom
        conservation at the final step).
  \item If the candidate passes all checks, the step is committed and
        added to the trajectory history. If not, the supervisor attempts
        to repair the violations using the correction strategies described
        in Section~\ref{sec:corrections}.
  \item If neither the original candidate nor any corrected version
        satisfies the constraints after the allowed number of retries,
        the supervisor backtracks: it reverts to the most recent valid
        state in the history and re-attempts from one timestep earlier,
        effectively undoing the failed step.
  \item The process terminates when $t = 0$ is reached or when the
        maximum backtrack budget is exhausted.
\end{enumerate}

The three possible outcomes for each timestep are:

\begin{itemize}[nosep]
  \item Commit: The denoised state passes all applicable checks on the
        first attempt. The step is accepted and appended to the
        trajectory history.
  \item Correct: The initial state fails a check, but one of the
        correction strategies (Section~\ref{sec:corrections}) repairs the
        violation. The corrected state replaces the original proposal.
  \item Backtrack: Neither the original proposal nor any corrected
        version passes. The supervisor reverts to the most recent valid
        state in the history and re-attempts from one timestep earlier.
\end{itemize}

The supervisor is parameterized by \texttt{max\_retries} (how many
correction attempts per timestep, default 3) and
\texttt{max\_backtracks} (global backtrack budget, default 10). These
ensure the loop terminates even in adversarial cases.

\subsection{Correction Strategies}
\label{sec:corrections}

When a denoising step produces an invalid molecular state, the
supervisor applies up to three targeted repair heuristics before
resorting to backtracking:

\begin{table}[H]
\centering
\caption{Correction strategies applied by the supervisor.}
\label{tab:corrections}
\begin{tabular}{@{}llp{6cm}@{}}
\toprule
Strategy & Targets & Mechanism \\
\midrule
\textsc{FixValency} & Overvalenced atoms &
  For each atom whose total bonds exceed its effective valency, reduce
  the explicit bond count (and then implicit hydrogen count if necessary)
  until the constraint is satisfied. \\[4pt]
\textsc{FixCharge} & Charge imbalance &
  Compute the charge deficit between target and current. Distribute
  $\pm 1$ charge adjustments across atoms, prioritizing lighter atoms
  first (lightest-first heuristic). \\[4pt]
\textsc{FixMass} & Mass imbalance &
  Compute the mass deficit. Add or remove implicit hydrogens
  (mass $\approx 1.008$\,u each) across atoms to bring the total mass
  within tolerance, respecting each atom's remaining bonding capacity. \\
\bottomrule
\end{tabular}
\end{table}

These heuristics are intentionally simple: they provide a best-effort
repair that is sufficient for the small molecules tested here. More
sophisticated strategies, such as asking Z3 to propose minimal edits,
are discussed in Section~\ref{sec:future}.

\subsection{Constraint-Aware Training}
\label{sec:training}

When the diffusion model is trained only on a standard denoising
objective, the supervisor must do all the work at inference time. To
reduce the frequency of constraint violations, ChemCS augments the
training loss with two differentiable penalty terms that encourage the
model to produce chemically plausible outputs.

The total training loss is a weighted sum of four components:
\begin{equation}
\mathcal{L} = w_1 \cdot \mathcal{L}_{\text{denoise}}
  + w_2 \cdot \mathcal{L}_{\text{bond}}
  + w_3 \cdot \text{scale}(e) \cdot \mathcal{L}_{\text{val}}
  + w_4 \cdot \text{scale}(e) \cdot \mathcal{L}_{\text{elem}}
\label{eq:loss}
\end{equation}

\noindent
where $\text{scale}(e)$ is the curriculum scaling factor
(Equation~\ref{eq:curriculum}) that ramps the constraint penalties from
zero to their full weight over the course of training. Table~\ref{tab:loss}
describes each term in detail.

\begin{table}[H]
\centering
\caption{Components of the training loss function.}
\label{tab:loss}
\begin{tabular}{@{}lp{7.5cm}r@{}}
\toprule
Term & Description & Default weight \\
\midrule
$\mathcal{L}_{\text{denoise}}$ &
  Mean squared error between the predicted atom features $\hat{x}_0$ and
  the clean target $x_0$. This is the standard denoising objective. &
  $w_1 = 1.0$ \\[4pt]
$\mathcal{L}_{\text{bond}}$ &
  Cross-entropy loss between the predicted bond-order logits
  (4 classes: 0, 1, 2, 3) and the ground-truth adjacency matrix. &
  $w_2 = 1.0$ \\[4pt]
$\mathcal{L}_{\text{val}}$ &
  Soft valency penalty:
  $\text{mean}\bigl(\text{ReLU}(\hat{b}_i - v_i)\bigr)$, where
  $\hat{b}_i$ is the expected total bond load (computed from softmax bond
  probabilities) and $v_i$ is the expected allowed valency (computed from
  softmax element probabilities). Fully differentiable. &
  $w_3 = 0.5$ \\[4pt]
$\mathcal{L}_{\text{elem}}$ &
  MSE between predicted and target element-count distributions, computed
  from the softmax of the element logits. Encourages the model to
  preserve the correct number of each element. &
  $w_4 = 0.3$ \\
\bottomrule
\end{tabular}
\end{table}

\subsubsection{Curriculum Schedule}

Introducing the constraint penalties ($\mathcal{L}_{\text{val}}$ and
$\mathcal{L}_{\text{elem}}$) at full strength from the start of training
destabilizes learning, because the model has not yet learned basic
denoising and the penalty gradients dominate the loss landscape.

ChemCS uses a curriculum schedule that sets the constraint penalty
weights to zero during a configurable warmup phase and then ramps them
linearly to their full values over the remaining epochs:
\begin{equation}
\text{scale}(e) = \begin{cases}
  0 & \text{if } e / E < f_{\text{warmup}} \\[2pt]
  \displaystyle\frac{e/E - f_{\text{warmup}}}{1 - f_{\text{warmup}}}
    & \text{otherwise}
\end{cases}
\label{eq:curriculum}
\end{equation}
where $e$ is the current epoch, $E$ is the total number of epochs, and
$f_{\text{warmup}}$ is the warmup fraction (default 25\%). The effective
penalty weights at epoch $e$ are $w_3 \cdot \text{scale}(e)$ and
$w_4 \cdot \text{scale}(e)$.

This allows the model to first learn how to denoise atom features and
bond orders from the standard objectives, and only then begin
incorporating chemical constraints once the representation is
sufficiently stable.

\subsubsection{Training Pipeline}

The PyTorch training pipeline mirrors the NumPy model architecture
(two-layer GNN with identical layer types) but adds a time embedding:
the fractional timestep $t/T \in [0, 1]$ is projected through a
two-layer MLP with SiLU activation and added to the atom features after
the input projection. This gives the model information about where it
is in the denoising trajectory.

Training proceeds for a configurable number of epochs (default 100).
At each epoch, for each training molecule, a random timestep is sampled,
noise is added to create the corrupted input, and the model predicts the
clean state. The loss (Equation~\ref{eq:loss}) is backpropagated with
gradient clipping (max norm 1.0) and Adam optimization (learning
rate $10^{-3}$).

After training, the learned weights are exported to the NumPy model via
\texttt{export\_to\_numpy()}, which copies all weight matrices and bias
vectors. This allows the supervisor loop to use trained weights without
requiring PyTorch at inference time.

% ===================================================================
% 5  EXPERIMENTAL SETUP
% ===================================================================
\section{Experimental Setup}

\subsection{Reaction Benchmark}

Experiments use three representative chemical reactions that exercise
different aspects of the constraint system:

\begin{table}[H]
\centering
\caption{Benchmark reactions used for evaluation.}
\label{tab:reactions}
\begin{tabular}{@{}llcl@{}}
\toprule
Reaction & Type & Atoms & Key challenge \\
\midrule
CH$_3$Br + OH$^-$ $\to$ ? & SN2 nucleophilic substitution & 7 &
  Formal charge transfer ($-1$) \\
CH$_4$ + 2\,O$_2$ $\to$ ? & Combustion & 9 &
  Multiple product molecules \\
HCl + NaOH $\to$ ? & Acid-base neutralization & 4 &
  Mass balance with light atoms \\
\bottomrule
\end{tabular}
\end{table}

\subsection{Experimental Conditions}

Each reaction is tested under two conditions at 50 random seeds:

\begin{itemize}[nosep]
  \item Unsupervised: Forward noise to $T{=}10$, then 10 reverse steps
        with no checks. The output is evaluated after completion.
  \item Supervised: Same model and noise, but wrapped in the supervisor
        loop with \texttt{max\_retries=2} and
        \texttt{max\_backtracks=3}.
\end{itemize}

Both conditions use an untrained GNN with hidden dimension 32, so that
the comparison isolates the effect of the supervisor from the effect of
learned weights.

\subsection{Metrics}

Three metrics are reported:

\begin{itemize}[nosep]
  \item Valency validity (\%): Fraction of runs in which every atom in
        the product respects its effective valency limit. This is the
        primary metric and can be measured mid-trajectory.
  \item Full conservation (\%): Fraction of runs in which all four
        axioms (mass, charge, atoms, valency) are satisfied. This is a
        much stricter criterion.
  \item Commit rate (\%): Fraction of denoising steps that passed on the
        first attempt without correction or backtracking. This measures
        how readily each step was accepted by the supervisor.
\end{itemize}

\subsection{Training Configuration}

For the training experiments, the model is trained on three small
molecules: water (H$_2$O), methane (CH$_4$), and ethanol (C$_2$H$_5$OH).
The training configuration is: hidden dimension 32, 2 layers, 200
epochs, learning rate $10^{-3}$, cosine noise schedule, and 25\% warmup
fraction for the curriculum. Training runs on CPU and completes in under
one minute.

% ===================================================================
% 6  RESULTS
% ===================================================================
\section{Results}

\subsection{Constraint Checker Validation}

The constraint checking module is validated through a comprehensive test
suite. Unit tests cover textbook chemical scenarios including:

\begin{itemize}[nosep]
  \item Valid reactions: SN2 (CH$_3$Br + OH$^-$ $\to$ CH$_3$OH + Br$^-$)
        and methane combustion (CH$_4$ + 2\,O$_2$ $\to$ CO$_2$ +
        2\,H$_2$O) both pass all four axioms.
  \item Valency violations: Carbon with 5 bonds, hydrogen with 2 bonds,
        and bromine with 2 bonds are all correctly flagged.
  \item Charge-adjusted valency: NH$_4^+$ (nitrogen with 4 bonds and
        $+1$ charge) is correctly accepted, since the $+1$ charge raises
        nitrogen's valency limit from 3 to 4.
  \item Conservation violations: Missing products (dropped Br$^-$),
        extra hydrogen atoms, and neutral Br replacing Br$^-$ are all
        detected.
  \item Edge cases: Empty product lists, single-atom molecules, and
        high-valency elements (sulfur with 6 bonds, phosphorus with
        5 bonds) are handled correctly.
\end{itemize}

All tests pass, covering the constraint checker, encoding/decoding,
noise schedule, reverse steps, correction heuristics, supervisor
integration, training loop, curriculum schedule, model export, and
dataset utilities.

\subsection{Supervised vs. Unsupervised Generation}

Figure~\ref{fig:valency} presents the main result: a comparison of
valency validity between supervised and unsupervised generation across
all three benchmark reactions.

\begin{figure}[H]
  \centering
  \includegraphics[width=0.85\textwidth]{figures/fig2_valency_comparison.png}
  \caption{Valency validity (\%) for supervised and unsupervised
           generation across three reactions, each tested with 50 random
           seeds using an untrained 32-dimensional GNN. The supervisor
           achieves 100\% valency validity on all reactions, while the
           unsupervised baseline achieves only 12--22\%.}
  \label{fig:valency}
\end{figure}

Table~\ref{tab:results} provides the complete numeric results:

\begin{table}[H]
\centering
\caption{Benchmark results: 50 random seeds per reaction, untrained
         GNN with hidden dimension 32, $T{=}10$ diffusion steps.}
\label{tab:results}
\begin{tabular}{@{}lcccc@{}}
\toprule
& \multicolumn{2}{c}{Valency validity} &
  \multicolumn{2}{c}{Full conservation} \\
\cmidrule(lr){2-3} \cmidrule(lr){4-5}
Reaction & Supervised & Unsupervised & Supervised & Unsupervised \\
\midrule
SN2         & 100\% & 12\% & 0\% & 0\% \\
Combustion  & 100\% & 12\% & 0\% & 0\% \\
Acid-base   & 100\% & 22\% & 0\% & 0\% \\
\midrule
Average     & 100\% & 15\% & 0\% & 0\% \\
\bottomrule
\end{tabular}
\end{table}

The key observations are:

\begin{enumerate}[nosep]
  \item The supervisor achieves 100\% valency validity across all three
        reactions and all 50 seeds for each. Without the supervisor, only
        about 15\% of runs produce valency-valid outputs.
  \item Full conservation remains at 0\% for both conditions. This is
        expected: the GNN has random weights and predicts essentially
        random element labels, so mass, charge, and atom-count
        conservation cannot be satisfied even with the supervisor's
        correction heuristics. Full conservation requires a trained model
        that produces chemically meaningful outputs.
\end{enumerate}

\subsection{Training Dynamics}

Figure~\ref{fig:training} shows the training loss over 200 epochs on
the three-molecule dataset.

\begin{figure}[H]
  \centering
  \includegraphics[width=0.85\textwidth]{figures/fig1_training_loss.png}
  \caption{Training loss components over 200 epochs on water, methane,
           and ethanol. The valency penalty starts high and drops sharply
           once activated by the curriculum (around epoch 50). Bond
           cross-entropy dominates the total loss throughout training.
           The denoising MSE converges quickly in the first 20 epochs.}
  \label{fig:training}
\end{figure}

The loss curve reveals several training phases:

\begin{itemize}[nosep]
  \item Epochs 1--50 (warmup): Only the denoising MSE and bond
        cross-entropy are active. The model learns basic feature
        reconstruction. The denoising MSE drops rapidly and stabilizes
        near zero.
  \item Epochs 50--100 (penalty ramp): The valency penalty is activated
        and initially spikes because the model has not yet learned to
        respect bond limits. It then decreases as the model adapts.
  \item Epochs 100--200 (steady state): All loss components stabilize.
        The bond cross-entropy remains the largest component, reflecting
        the difficulty of predicting the correct bond structure among
        $O(n^2)$ atom pairs.
\end{itemize}

\subsection{Curriculum Schedule Ablation}

Figure~\ref{fig:curriculum} compares three curriculum configurations to
illustrate the effect of the warmup fraction.

\begin{figure}[H]
  \centering
  \includegraphics[width=0.7\textwidth]{figures/fig3_curriculum.png}
  \caption{Curriculum penalty weight scale over 100 epochs for three
           warmup strategies. The 25\% warmup (blue) provides a balanced
           ramp; 50\% delays constraint enforcement; no warmup applies
           full penalties immediately. In practice, 25\% was found to
           give the most stable training convergence.}
  \label{fig:curriculum}
\end{figure}

The 25\% warmup was selected as the default after observing that using
no warmup caused loss spikes in early training (the constraint penalties
overwhelmed the denoising gradients before the model could learn basic
feature reconstruction), while 50\% warmup left too few epochs for the
penalties to take effect.

% ===================================================================
% 7  DISCUSSION
% ===================================================================
\section{Discussion}

The central finding is that constraint enforcement during generation
substantially improves local chemical validity, even when the underlying
neural network has not been trained. The supervisor achieves 100\%
valency validity with random weights, demonstrating that the benefit
comes from the logical checking and repair mechanism rather than from
the model itself.

This result supports the broader idea that logic and learning are
complementary: the neural network provides the generative capacity
(proposing plausible structures from noise), while the constraint system
provides the domain knowledge (rejecting or repairing implausible
configurations). Neither component alone is sufficient. A constraint
checker without a generator cannot propose novel molecules, and a
generator without constraints produces mostly invalid outputs.

An important nuance is the interaction between inference-time correction
and training. When the supervisor fixes the model's mistakes at
inference time, the model does not receive any signal about those
mistakes. This is why the constraint-aware training loss is important:
it gives the model gradient signal about chemical validity during
learning, so that over time the supervisor has fewer violations to
correct. The curriculum schedule is essential here because the denoising
and constraint objectives operate at different scales and the model
needs to learn basic feature reconstruction before it can meaningfully
respond to valency penalties.

The runtime overhead of the supervisor is moderate and could be reduced
with batched constraint checking or by running the Z3 solver only on
suspicious atoms (those with high predicted bond counts). For the small
molecules tested here the overhead is negligible in absolute terms.

% ===================================================================
% 8  LIMITATIONS AND FUTURE WORK
% ===================================================================
\section{Limitations and Future Work}
\label{sec:future}

\emph{Training data scale.} The training experiments use only three
molecules (water, methane, ethanol). A serious evaluation would require
training on a standard dataset such as QM9~\cite{ramakrishnan2014qm9}
or ZINC~\cite{irwin2012zinc}, parsed from SMILES strings via RDKit.

\emph{Molecular complexity.} The current system handles molecules up to
approximately 15 atoms. The bond head's pairwise computation is $O(n^2)$
in the number of atoms, which becomes expensive for larger molecules.
There are no rules for ring structures, stereochemistry, or aromaticity.

\emph{Full conservation.} With random GNN weights, mass and atom-count
conservation cannot be achieved because the model predicts arbitrary
element types. Trained weights and a larger dataset are prerequisites
for full-conservation validity.

\emph{Correction sophistication.} The current repair heuristics
(Table~\ref{tab:corrections}) are greedy and local: they fix one atom
at a time without considering the global molecular structure. A more
principled approach would be to ask the Z3 solver to find a minimal set
of edits that restores validity, treating the correction itself as a
constraint-satisfaction problem.

\emph{Benchmark scale.} The current benchmark uses 50 seeds across 3
reactions. Scaling to approximately 1,000 reactions across diverse
reaction types would provide stronger statistical evidence.

\emph{Comparison to existing methods.} This work does not compare
against established molecular generation baselines (e.g., DiGress,
GDSS, or MolGAN). Such comparisons would be necessary to position
ChemCS relative to the state of the art.

% ===================================================================
% 9  CONCLUSION
% ===================================================================
\section{Conclusion}

This paper presented ChemCS, a system that integrates constraint
satisfaction directly into the reverse-diffusion loop of a graph neural
network molecular generator. By checking chemical axioms (valency, mass
conservation, charge balance, and atom counts) at every denoising step,
the supervisor achieves 100\% valency validity on a benchmark of three
representative reactions, compared with approximately 15\% for the same
model without supervision. A companion training pipeline adds
differentiable constraint penalties through a curriculum schedule,
providing gradient signal that complements the inference-time checks.

The architecture demonstrates that combining symbolic reasoning (via Z3)
with neural generation (via a GNN denoiser) can enforce domain
constraints without sacrificing the flexibility of the generative model.
The code, figures, and tests are available in the project repository.
The NumPy inference path runs on CPU; PyTorch is required only for
training.

% ===================================================================
% REFERENCES
% ===================================================================
\begin{thebibliography}{9}

\bibitem{hoogeboom2022equivariant}
E.~Hoogeboom, V.~G.~Satorras, C.~Vignac, and M.~Welling,
``Equivariant diffusion for molecule generation in 3D,''
in \textit{Proc.\ International Conference on Machine Learning (ICML)},
vol.~162, pp.~8867--8887, 2022.

\bibitem{vignac2023digress}
C.~Vignac, I.~Krawczuk, A.~Siraudin, B.~Wang, V.~Cevher, and P.~Frossard,
``DiGress: Discrete denoising diffusion for graph generation,''
in \textit{Proc.\ International Conference on Learning Representations
(ICLR)}, 2023.

\bibitem{nichol2021improved}
A.~Q.~Nichol and P.~Dhariwal,
``Improved denoising diffusion probabilistic models,''
in \textit{Proc.\ International Conference on Machine Learning (ICML)},
vol.~139, pp.~8162--8171, 2021.

\bibitem{demoura2008z3}
L.~de~Moura and N.~Bj{\o}rner,
``Z3: An efficient SMT solver,''
in \textit{Proc.\ Tools and Algorithms for the Construction and Analysis
of Systems (TACAS)}, Lecture Notes in Computer Science, vol.~4963,
pp.~337--340, 2008.

\bibitem{ramakrishnan2014qm9}
R.~Ramakrishnan, P.~O.~Dral, M.~Rupp, and O.~A.~von Lilienfeld,
``Quantum chemistry structures and properties of 134 kilo molecules,''
\textit{Scientific Data}, vol.~1, no.~140022, 2014.

\bibitem{irwin2012zinc}
J.~J.~Irwin, T.~Sterling, M.~M.~Mysinger, E.~S.~Bolstad, and R.~G.~Coleman,
``ZINC: A free tool to discover chemistry for biology,''
\textit{Journal of Chemical Information and Modeling}, vol.~52, no.~7,
pp.~1757--1768, 2012.

\end{thebibliography}

\end{document}
